\documentclass[10.5pt,a4paper]{article}
\usepackage[a4paper,margin=2.2cm,top=2.4cm,bottom=2.4cm]{geometry}
\usepackage{fontspec}
\usepackage{amsmath,amssymb,mathtools,bm}
\setmainfont{Noto Serif CJK SC}[BoldFont={Noto Serif CJK SC Bold}]
\setsansfont{Noto Sans CJK SC}[BoldFont={Noto Sans CJK SC Bold}]
\XeTeXlinebreaklocale "zh"
\XeTeXlinebreakskip = 0pt plus 1pt minus 0.1pt
\linespread{1.18}
\usepackage{xcolor}
\definecolor{accent}{HTML}{1F5FA8}
\definecolor{quiet}{HTML}{5A5A5A}
\definecolor{tint}{HTML}{EEF3FA}
\usepackage{titlesec}
\titleformat{\section}{\sffamily\bfseries\Large\color{accent}}{\thesection}{0.8em}{}
\titleformat{\subsection}{\sffamily\bfseries\large}{\thesubsection}{0.6em}{}
\titleformat{\subsubsection}{\sffamily\bfseries\normalsize}{\thesubsubsection}{0.5em}{}
\titlespacing*{\section}{0pt}{1.6em}{0.7em}
\titlespacing*{\subsection}{0pt}{1.2em}{0.5em}
\usepackage{enumitem}
\setlist{itemsep=0.25em,topsep=0.35em,leftmargin=1.6em}
\usepackage{array,booktabs,tabularx,xltabular,makecell}
\renewcommand{\arraystretch}{1.25}
\newcolumntype{L}{>{\raggedright\arraybackslash}X}
\newcolumntype{P}[1]{>{\raggedright\arraybackslash}p{#1}}
\usepackage{caption}
\captionsetup{font=small,labelfont=bf,labelsep=quad}
\renewcommand{\figurename}{图}
\renewcommand{\tablename}{表}
\renewcommand{\contentsname}{目录}
\usepackage{graphicx}\usepackage{float}
\usepackage{tikz}
\usetikzlibrary{positioning,arrows.meta,calc,fit}
\usepackage{fancyhdr}
\pagestyle{fancy}\fancyhf{}
\fancyhead[L]{\small\color{quiet}椭圆通道（第 3 版）：参数、强度场与定义的适定性}
\fancyhead[R]{\small\color{quiet}2026 年 10 月 16 日}
\fancyfoot[C]{\small\thepage}
\renewcommand{\headrulewidth}{0.4pt}
\usepackage[colorlinks=true,linkcolor=accent,citecolor=accent,urlcolor=accent]{hyperref}
\newcommand{\AFH}{AFRC$_{\mathrm H}$}
\newcommand{\AFE}{AFRC$_{\mathrm E}$}
\newcommand{\lnJ}{$\ln J$}
\newcommand{\keyp}[1]{\textbf{\sffamily #1}}
\newenvironment{finding}{\begin{center}\begin{minipage}{0.94\linewidth}\setlength{\parskip}{0.3em}\color{black}}{\end{minipage}\end{center}}
\setlength{\parskip}{0.35em}
\setlength{\parindent}{0pt}
\tolerance=2000 \emergencystretch=3em
\usepackage{amsthm}
\theoremstyle{definition}
\newtheorem{definition}{定义}[section]
\theoremstyle{plain}
\newtheorem{proposition}[definition]{命题}
\newtheorem{corollary}[definition]{推论}
\theoremstyle{remark}
\newtheorem{remark}[definition]{注}
\renewcommand{\proofname}{证明}
\newcommand{\status}[1]{\hfill{\small\color{accent}\textsf{[#1]}}}
\newcommand{\re}{\operatorname{Re}}
\newcommand{\im}{\operatorname{Im}}
\begin{document}
\begin{center}
{\sffamily\bfseries\LARGE 椭圆通道（第 3 版）：参数、强度场与定义的适定性}\\[0.5em]
{\sffamily\large 平滑与阈值扫描；同步簇层级作为强度场；LAVD 在涡片上是否定义良好}\\[0.8em]
{\color{quiet}2026 年 10 月 16 日\quad 接第 2 版}
\end{center}
\vspace{0.3em}
\noindent\fcolorbox{accent!40}{tint}{\begin{minipage}{0.96\linewidth}\vspace{0.4em}
\keyp{结论}
\begin{itemize}[leftmargin=1.4em,itemsep=0.1em]
\item \keyp{参数。}
  \begin{itemize}[leftmargin=1.2em]
  \item 平滑核宽度实测为 $w(n)\approx0.55\sqrt n\,h$。尺寸上限怎么收紧，9 个参照涡 × 3 个噪声种子的平均噪声 IoU 都只变化 $\le0.02$（$\delta=0.1$ 时 0.626–0.629，逐涡最优为 0.681）。最优平滑次数不随涡的尺寸变化，剩余差距由各涡的信噪比决定。第 3 版采用 $w\le r_{\rm eq}/6$。
  \item $\delta$ 没有唯一最优值，建议报告相干半径随 $\delta$ 的曲线。
  \item 持续转动门：单步占比的余量窄（真核 $\le0.23$，Sod 阈值以下的伪核 $\ge0.27$）。按比例的参与比 $n_{\rm eff}/(S-1)$ 会误杀激波–气泡中窗口内才形成的涡（0.25），不采用。第 3 版取 $J_{\max}=0.35$。
  \end{itemize}
\item \keyp{强度场的网络表达。}同步文献里“强度”就是频率失谐，即 $\varphi/T$，与 LAVD 的被积量相同；同步框架给不出另一种强度，但频率簇的层级（$\varphi$ 的合并树）给出另一种参照：
  \begin{itemize}[leftmargin=1.2em]
  \item 背景自然为零。
  \item 对集合参照常数严格不变：双环流去掉部分示踪点后 $\varphi$ 平移 0.26 rad，核心 $P$ 不变到小数点后 4 位。
  \item 代价：要求簇的环境落在观测区域内（Lamb–Oseen 小圆盘中核心 $P$ 由 18.3 降到 11.7）；噪声下在激波–涡上比 $\varphi$ 差。
  \end{itemize}
\item \keyp{LAVD 会不会定义不明：会，问题在涡片上。}网络的 $\varphi$ 只用位置，但它与 LAVD 一样是逐点的转角场，在涡片（滑移线、带剪切的接触面）上没有分辨率无关的极限：层厚远小于 $h$ 时，节点 $\max|\varphi|$ 随 $h$ 从 8.2 涨到 23.1，随帧数从 13.1 涨到 21.7，连区域积分也不收敛。物质环路的环量则严格守恒、与采样无关。由此给出基于环量的强度：在相干边界上取环量，以外侧环带为参照。它只用速度的线积分（速度由位置差分得到），客观性残差 0.2\%–0.7\%，在光滑流中与 $\varphi$ 一致到 2\%–11\%。
\end{itemize}\vspace{0.2em}\end{minipage}}

\tableofcontents

\section{平滑、$\delta$ 与 $J_{\max}$}
\subsection{平滑核宽度与尺寸上限}
\begin{table}[H]\centering\small\caption{图平滑核的均方根宽度 $/h$（在 20 个节点上对 $\delta$ 函数平滑后取中位数）。}\begin{tabular}{lccccccc}\toprule 算例 & $n=2$ & $n=3$ & $n=4$ & $n=6$ & $n=8$ & $n=12$ & $n=16$\\\midrule
Lamb–Oseen & 0.79 & 0.97 & 1.12 & 1.36 & 1.57 & 1.91 & 2.21\\
双环流 & 0.81 & 0.99 & 1.13 & 1.38 & 1.59 & 1.93 & 2.22\\
激波–涡 & 0.79 & 0.96 & 1.10 & 1.35 & 1.55 & 1.89 & 2.18\\
激波–气泡 & 0.78 & 0.95 & 1.09 & 1.34 & 1.55 & 1.90 & 2.19\\
$0.55\sqrt n$ & 0.78 & 0.95 & 1.10 & 1.35 & 1.56 & 1.91 & 2.20\\\bottomrule\end{tabular}\end{table}
\begin{table}[H]\centering\small\caption{9 个参照涡、3 个噪声种子（激波–气泡为 2 个）下，噪声边界与参照边界 IoU 的平均值。}\begin{tabular}{lcc}\toprule 平滑规则 & $\delta=0.1$ & $\delta=0.2$\\\midrule
固定 2 次 & 0.523 & 0.586\\
固定 8 次（噪声规则） & 0.628 & 0.656\\
第 2 版上限 $4(r/3h)^2$ & 0.629 & 0.658\\
$w\le r/3$ & 0.628 & 0.656\\
$w\le r/4$ & 0.626 & 0.666\\
$w\le r/6$（第 3 版） & 0.628 & 0.675\\
逐涡最优（不可实现） & 0.681 & 0.755\\
\bottomrule\end{tabular}\end{table}
\begin{table}[H]\centering\small\caption{各涡在不同平滑次数下的噪声 IoU（$\delta=0.1$，种子平均）。}\resizebox{\linewidth}{!}{\begin{tabular}{llccccccccc}\toprule 算例 & 涡 & $r_{\rm eq}/h$ & 最优 $n$ & 2 & 3 & 4 & 6 & 8 & 12 & 16\\\midrule
Lamb–Oseen & (-3.04, 0.04) & 3.2 & 16 & 0.79 & 0.86 & 0.87 & 0.80 & 0.87 & 0.85 & 0.87\\
Lamb–Oseen & (3.00, -0.07) & 3.4 & 3 & 0.93 & 0.95 & 0.94 & 0.93 & 0.93 & 0.93 & 0.91\\
双环流 & (1.50, 0.37) & 18.4 & 8 & 0.56 & 0.72 & 0.73 & 0.75 & 0.79 & 0.79 & 0.79\\
双环流 & (0.75, 0.73) & 17.2 & 16 & 0.70 & 0.85 & 0.87 & 0.90 & 0.90 & 0.90 & 0.92\\
激波–涡 & (0.70, 0.50) & 20.8 & 16 & 0.45 & 0.46 & 0.60 & 0.61 & 0.66 & 0.66 & 0.67\\
激波–涡 & (0.85, 0.52) & 5.4 & 12 & 0.31 & 0.32 & 0.38 & 0.38 & 0.40 & 0.43 & 0.38\\
激波–涡 & (0.55, 0.48) & 17.5 & 16 & 0.00 & 0.00 & 0.04 & 0.17 & 0.16 & 0.15 & 0.39\\
激波–涡 & (0.95, 0.65) & 6.0 & 3 & 0.41 & 0.55 & 0.42 & 0.50 & 0.46 & 0.44 & 0.41\\
激波–气泡 & (1.00, 0.04) & 5.5 & 2 & 0.56 & 0.56 & 0.54 & 0.52 & 0.48 & 0.46 & 0.44\\
\bottomrule\end{tabular}}\end{table}
各种尺寸上限之间相差不超过 0.02，与固定 8 次几乎相同。最优平滑次数不随 $r_{\rm eq}/h$ 变化：
\begin{itemize}
\item Lamb–Oseen 的涡半径约 3$h$，从 3 次到 16 次都差不多；
\item 激波–涡半径 6$h$ 的涡最优为 3 次；
\item 激波–气泡半径 5.5$h$ 的涡最优为 2 次。
\end{itemize}
所以尺寸不是控制量，剩下的约 0.05 的差距由各涡的信噪比决定。第 3 版取 $w(n)\le r_{\rm eq}/6$，它在 $\delta=0.2$ 时最好（0.675），在 $\delta=0.1$ 时持平。

\subsection{$\delta$ 的敏感性}
\begin{table}[H]\centering\small\caption{每格为“相干半径 $/h$ / 与参照的干净 IoU / 噪声 IoU（8 次平滑）”。}\resizebox{\linewidth}{!}{\begin{tabular}{llccccc}\toprule 算例 & 涡 & $\delta=0.05$ & $\delta=0.1$ & $\delta=0.15$ & $\delta=0.2$ & $\delta=0.3$\\\midrule
Lamb–Oseen & (-3.04, 0.04) & 2.9 / 1.00 / 0.71 & 3.2 / 0.97 / 0.87 & 3.6 / 1.00 / 0.91 & 3.9 / 0.91 / 0.86 & 4.2 / 1.00 / 0.76\\
Lamb–Oseen & (3.00, -0.07) & 3.3 / 1.00 / 0.93 & 3.4 / 1.00 / 0.93 & 3.6 / 0.95 / 0.89 & 4.4 / 0.97 / 0.70 & 4.7 / 0.96 / 0.88\\
双环流 & (1.50, 0.37) & 15.2 / 0.58 / 0.67 & 18.4 / 0.86 / 0.79 & 19.5 / 0.94 / 0.88 & 20.1 / 0.94 / 0.93 & 21.3 / 0.95 / 0.90\\
双环流 & (0.75, 0.73) & 15.9 / 0.99 / 0.95 & 17.2 / 0.92 / 0.90 & 17.8 / 1.00 / 0.89 & 20.1 / 0.82 / 0.74 & 22.9 / 0.95 / 0.72\\
激波–涡 & (0.70, 0.50) & 17.6 / 0.94 / 0.66 & 20.8 / 0.86 / 0.66 & 20.8 / 1.00 / 0.75 & 20.8 / 1.00 / 0.92 & 20.8 / 1.00 / 0.98\\
激波–涡 & (0.85, 0.52) & 5.4 / 0.84 / 0.37 & 5.4 / 0.84 / 0.40 & 5.4 / 0.84 / 0.53 & 5.4 / 0.84 / 0.57 & 5.4 / 0.84 / 0.56\\
激波–涡 & (0.55, 0.48) & 17.5 / 0.78 / 0.16 & 17.5 / 0.97 / 0.16 & 17.5 / 0.97 / 0.16 & 17.5 / 0.97 / 0.16 & 17.5 / 0.97 / 0.52\\
激波–涡 & (0.95, 0.65) & 6.0 / 0.91 / 0.46 & 6.0 / 0.91 / 0.46 & 6.0 / 0.91 / 0.46 & 6.0 / 0.91 / 0.46 & 6.0 / 0.91 / 0.46\\
激波–气泡 & (1.00, 0.04) & 4.5 / 1.00 / 0.70 & 5.5 / 0.90 / 0.48 & 5.7 / 0.93 / 0.50 & 5.7 / 0.93 / 0.58 & 5.8 / 0.93 / 0.81\\
\bottomrule\end{tabular}}\end{table}
\begin{itemize}
\item 相干半径随 $\delta$ 增大：Lamb–Oseen 从 2.9$h$ 到 4.7$h$，双环流从 15$h$ 到 23$h$。激波–涡主涡从 $\delta=0.1$ 起停在 20.8$h$，说明它的边界由别的结构（相邻涡或等值线终止）决定。
\item 干净 IoU 几乎都在 0.84 以上；噪声 IoU 在 $\delta=0.15$–0.2 时平均略高。
\end{itemize}
与 RCLV 研究中按参数划分相干等级的做法类似 [7]，建议报告半径–$\delta$ 曲线，不固定单一的 $\delta$。

\subsection{持续转动门：单步占比与参与比}
\begin{table}[H]\centering\small\caption{干净数据上前 40 个核的单步占比与参与比 $n_{\rm eff}=(\sum|x|)^2/\sum x^2$。}\resizebox{\linewidth}{!}{\begin{tabular}{lccc}\toprule 算例 & $S-1$ & 显著核（持续性 $>0.5$）：单步占比 / $n_{\rm eff}$ & 其余核：单步占比 / $n_{\rm eff}$\\\midrule
Sod & 40 & — / — & 0.27–10.36 / 1.12–10.79\\
激波–涡 & 30 & 0.04–0.06 / 25.82–29.22 & 0.04–0.29 / 10.11–29.41\\
Lamb–Oseen & 40 & 0.03–0.03 / 39.38–40.00 & 0.03–0.03 / 38.46–39.99\\
双环流 & 50 & 0.03–0.06 / 30.87–47.69 & 0.03–0.63 / 15.71–45.95\\
激波–气泡 & 70 & 0.02–0.23 / 17.50–57.86 & 0.02–0.61 / 16.76–57.96\\
\bottomrule\end{tabular}}\end{table}
\begin{itemize}
\item \keyp{单步占比}：真核不超过 0.23；但 Sod 低于显著性阈值的伪核可以低到 0.27（两三次激波与接触间断事件）。余量很窄，主要靠持续性阈值把它们挡住。
\item \keyp{参与比}：Sod 伪核不超过 10.8 步。但激波–气泡显著核只有 17.5–57.9/70 步，因为这些涡在窗口内才形成。按比例的门会误杀它们，所以不采用。
\item \keyp{噪声下两者都失效}：噪声把转角分散到每一步，但显著性门已经挡住了这类噪声核。
\end{itemize}
第 3 版取 $J_{\max}=0.35$。在超过持续性阈值的核中，Sod 唯一的假核为 1.00，真核不超过 0.23。

\section{强度场的网络表达：同步簇与合并树}
\subsection{文献}
\begin{itemize}
\item \keyp{高阶拓扑同步} [1–3]：相位可以定义在边与三角形上，但边上的同步只有投影到无旋/无散分量之后才看得出来 [1]。这与本工作的发现直接对应：点对转角混有应变，三角形转角就是边信号的旋转部分。节点与边信号的耦合 [2]、Dirac 同步 [3] 对应我们的 T5 Dirac 结构。
\item \keyp{簇同步} [4]：从网络对称性预测同步簇，前提是动力学由网络本身产生。我们的图只是观测图，对称性分析不适用；能借用的是“频率簇”这个概念。
\item \keyp{Chimera 诊断} [5]：用平均相速度剖面与离散 Laplacian（局部曲率）度量空间相干性，阈值取最大曲率的 1\%。这与 \AFE 的曲率超额同属 Laplacian 型对比，零点也同样是相对的。
\item \keyp{轨迹时间平均与聚类}：遍历商 [6]、谱聚类 [8]。
\end{itemize}
在同步语言里，涡的“强度”就是频率失谐 $\Omega_i-\bar\Omega=\varphi_i/T$，与 LAVD 的被积量相同。所以同步框架不会给出另一种强度，但会给出另一种参照：用频率簇的层级代替全局均值。

\subsection{合并树强度场}
\begin{definition}[频率簇层级强度]
对 $\varphi$ 做上水平集过滤。节点的 $P^+_i$ 是它高出所在簇“终止水平”的高度，终止水平取以下两者中先发生的：并入峰值更高的簇；触到观测物质的边缘。$P=P^+(\varphi)-P^+(-\varphi)$。$P$ 就是 $\varphi$ 的 0 维持续同调，对 $\varphi\to\varphi+c$ 严格不变。
\end{definition}
\begin{table}[H]\centering\small\caption{合并树强度场与参照（同一定义作用于参照 $\varphi$）的一致性。噪声下 $P$ 用 8 次平滑。}\resizebox{\linewidth}{!}{\begin{tabular}{lcccccc}\toprule 算例 & 涡区占比 & Spearman $P$ 干净 / 噪声 & Spearman $\varphi$ 干净 / 噪声 & Pearson $P$ / $\varphi$（噪声） & 背景 $|\varphi|$/核 & 背景 $|P|$/核\\\midrule
Lamb–Oseen & 0.17 & 1.000 / 0.999 & 0.999 / 0.991 & 0.999 / 0.997 & 0.050 & 0.000\\
双环流 & 0.45 & 1.000 / 1.000 & 1.000 / 0.997 & 0.999 / 0.999 & 0.187 & 0.000\\
激波–涡 & 0.34 & 0.972 / 0.597 & 0.998 / 0.859 & 0.683 / 0.992 & 0.035 & 0.000\\
激波–气泡 & 0.11 & 0.992 / 0.950 & 0.996 / 0.978 & 0.958 / 0.987 & 0.027 & 0.000\\
\bottomrule\end{tabular}}\end{table}
\begin{table}[H]\centering\small\caption{去掉部分示踪点（Lamb–Oseen 只保留 $r_0<4.5$，双环流只保留 $x_0<1.85$）。}\resizebox{\linewidth}{!}{\begin{tabular}{lccccc}\toprule 算例 & 保留示踪点 & $\varphi$ 平移（中位数 / 四分位距） & 核心 $\varphi$ 全集 / 子集 & 核心 $P$ 全集 / 子集\\\midrule
Lamb–Oseen & 3839 & -0.928 / 7e-16 & 17.748 / 16.820 & 18.277 / 11.702\\
双环流 & 11840 & +0.261 / 1e-15 & -9.774 / -9.513 & -7.121 / -7.121\\
\bottomrule\end{tabular}}\end{table}
\begin{itemize}
\item $\varphi$ 平移的是一个精确常数（四分位距 $10^{-15}$），证实 $\varphi$ / LAVD 依赖播撒的集合。这也是论文结论中写到的局限。
\item 双环流涡核的 $P$ 完全不变。Lamb–Oseen 小圆盘切到了涡的外围旋流，簇先触到边缘，$P$ 从 18.3 降到 11.7：$P$ 是局部量，前提是簇的环境要在观测区域内。
\item 两个被否定的变体：
  \begin{itemize}[leftmargin=1.2em]
  \item “覆盖半个区域的簇算背景”在双环流上失效，因为两个涡填满了区域。
  \item 以涡自身相干边界为参照的强度在理论上严格不变（实测残差 0.5\%–4\% 来自等值线扫描的离散化），但固体旋转涡会得到约 0：它量的是核心相对自身边缘的转动（Lamb–Oseen 为 2.9 rad，合并树为 18.3），物理上不对。
  \end{itemize}
\item 噪声下：Lamb–Oseen、双环流、激波–气泡上 $P$ 与 $\varphi$ 相当；激波–涡上 $P$ 明显更差（Spearman 0.60 对 0.86），因为噪声会移动合并鞍点的水平。
\end{itemize}

\section{LAVD 会不会定义不明}
\subsection{问题在哪里}
\begin{itemize}
\item 网络本身不用速度梯度，也不用涡量：$\varphi$ 来自三角形步映射的 $\arg p$，只用位置。LAVD 只以两种身份出现：一是 $\varphi$ 在光滑流中的连续极限；二是验证时用的 DNS 涡量参照。
\item 但目标本身有问题。在 BV 意义下，滑移线和带剪切的接触面上，涡量是集中在曲线上的奇异测度：
  \begin{itemize}[leftmargin=1.2em]
  \item 逐点的 $\int|\omega-\bar\omega|\,dt$ 在片上为无穷，片附近也没有分辨率无关的极限；
  \item 任何逐点的转角场（包括 $\varphi$ 和 $P$）在片上都会继承这一问题；
  \item 激波生成的涡量只是有界跳跃（Crocco），逐点几乎处处有定义，问题不大；
  \item 参照 $\bar\omega$ 依赖集合（第 2.2 节）。
  \end{itemize}
\item 积分量不受影响：涡量测度在区域上的积分，即物质环路的环量，是良定义的（离散 Kelvin 定理，统一理论 T4）。
\end{itemize}

\subsection{涡片检验}
$u=(U\tanh(y/\delta_s),0)$，$U=0.5$，$T=1$。$\delta_s=0.1$ 为可分辨的剪切层，$\delta_s=5\times10^{-4}$ 为远小于 $h$ 的涡片。环量取约 $0.2\times0.2$ 的物质方框边界上的示踪点环路，用梯形规则，速度由位置的时间差分得到。
\begin{table}[H]\centering\small\caption{剪切层与涡片上的逐点转角、区域积分与环路环量。}\resizebox{\linewidth}{!}{\begin{tabular}{cccccccc}\toprule 层厚 $\delta_s$ & $h$ & 帧数 & 节点 $\max|\varphi|$ & 层内 $|\varphi|$ 中位数 & 区域积分 $\int\varphi\,dA$ & $\Gamma(t_0)$ & $\Gamma(t_{\rm end})$\\\midrule
0.1 & 0.04 & 41 & 1.58 & 1.48 & -0.041 & -0.0878 & -0.0878\\
0.1 & 0.02 & 41 & 1.66 & 1.65 & -0.045 & -0.1203 & -0.1203\\
0.1 & 0.01 & 41 & 1.67 & 1.67 & -0.043 & -0.1344 & -0.1344\\
0.1 & 0.02 & 161 & 1.66 & 1.65 & -0.045 & -0.1203 & -0.1203\\
0.0005 & 0.04 & 41 & 8.20 & 5.31 & -0.082 & -0.1549 & -0.1549\\
0.0005 & 0.02 & 41 & 13.15 & 10.51 & -0.063 & -0.1781 & -0.1781\\
0.0005 & 0.01 & 41 & 23.14 & 15.08 & -0.044 & -0.1836 & -0.1836\\
0.0005 & 0.02 & 161 & 21.72 & 14.44 & -0.081 & -0.1781 & -0.1781\\
\bottomrule\end{tabular}}\end{table}
\begin{itemize}
\item \keyp{可分辨层}：$\max|\varphi|$ 收敛到精确值 $\tfrac12(\omega-\bar\omega)T=1.67$，与帧数无关。
\item \keyp{涡片}：
  \begin{itemize}[leftmargin=1.2em]
  \item 节点转角随 $h$ 和帧间隔都发散：$h$ 从 0.04 到 0.01 时由 8.2 涨到 23.1；帧数从 41 到 161 时由 13.1 涨到 21.7。
  \item 区域积分同样不收敛。原因是剪切三角形的单步转角是 $\arctan(\gamma/2)$，会饱和。
  \item 环量在每种设置下都严格守恒，且与帧数无关（$-0.1781$ 对 $-0.1781$）。
  \end{itemize}
\end{itemize}

\subsection{基于环量的强度}
\begin{definition}[环量强度]
$C$ 为相干边界（周长判据），$R$ 为 $C$ 与外侧一条等值线之间的环带，取面积为 $C$ 的 1.8–3 倍。定义
\begin{equation}
S_\Gamma=\frac12\,T\,\Big\langle \frac{\Gamma_C(t)}{A_C(t)}-\frac{\Gamma_{C\cup R}(t)-\Gamma_C(t)}{A_{C\cup R}(t)-A_C(t)}\Big\rangle_t .
\end{equation}
\end{definition}
\begin{itemize}
\item 它只用物质环路上速度的线积分，速度由示踪点位置差分得到。在网络上就是速度 1-上链 $f$ 与边界 1-链 $\partial C$ 的配对，与内部剖分无关（统一理论命题 3）。
\item 参考系转动给两项各加 $2\Omega$，相互抵消，所以客观。
\item 参照是外侧环带，不是全局均值，所以不依赖集合。
\end{itemize}
\begin{table}[H]\centering\small\caption{环量强度与 $\varphi$ 型强度（$C$ 内平均减环带平均）的比较。旋转参考系为窗口内转 0.6 rad 并平移。}\resizebox{\linewidth}{!}{\begin{tabular}{llccccc}\toprule 算例 & 涡 & 面积 & $S_\Gamma$（环量） & $S_\Gamma$（旋转参考系） & $S_\varphi$（内减环带） & $\varphi_{\rm core}$\\\midrule
Lamb–Oseen & (-3.04, -0.04) & 0.5226 & 2.320 & 2.304 & 2.605 & 17.75\\
Lamb–Oseen & (3.03, 0.07) & 0.6012 & 2.584 & 2.567 & 2.803 & 17.74\\
双环流 & (1.51, 0.38) & 0.1445 & 2.039 & 2.035 & 2.075 & 9.79\\
双环流 & (0.75, 0.73) & 0.1336 & -1.882 & -1.886 & 1.906 & -9.77\\
激波–涡 & (0.59, 0.60) & 0.0242 & -0.112 & -0.112 & 0.109 & -1.22\\
\bottomrule\end{tabular}}\end{table}
\begin{itemize}
\item 双环流两者一致到 2\%，Lamb–Oseen 一致到 9\%–11\%。偏差来自外侧环路处在剪切中，P1 像不够准。
\item 旋转参考系下的残差为 0.2\%–0.7\%，来自旋转后位置的有限差分。
\item 激波–涡只有第三个涡找到了合适的外侧环路（0.112 对 0.109），另两个在面积达到 1.8 倍之前就与相邻结构合并；激波–气泡一个也没有找到。外侧环路的选取规则还需要改进，例如取合并树的终止等值线。
\end{itemize}

\begin{finding}
\keyp{回答}\quad 用 LAVD 作为强度，确实会把涡量的逐点定义带进来，而在涡片上它没有良定义的极限。网络的 $\varphi$ 虽然只用位置，作为逐点场也有同样的问题。适定的做法是：
\begin{itemize}[leftmargin=1.2em]
\item \keyp{边界}：由 \AFE 导出的周长判据给出；跨越涡片的环路必然拉伸，自动不被判为相干。
\item \keyp{强度}：相干环路上的环量，以外侧环带为参照。
\item \keyp{$\varphi$ 的角色}：只用来提供候选等值线与核的位置，不再作为强度报告。
\end{itemize}
\end{finding}

\section{第 3 版六算例复测}
\begin{table}[H]\centering\small\caption{持续性核数与“涡”数（$J_{\max}=0.35$）。}\resizebox{\linewidth}{!}{\begin{tabular}{lcccccccc}\toprule 算例 & 参照核 & 旧法 干净 & 旧法 噪声（匹配） & \textbf{软阈值 噪声}（匹配） & 持续门剔除 干净/噪声 & 参照涡 & \textbf{涡 干净}（匹配） & \textbf{涡 噪声}（匹配）\\\midrule
双环流 & 3 & 3 & 4（3） & \textbf{3}（3） & 0/0 & 2 & \textbf{2}（2） & \textbf{2}（2） \\
Lamb–Oseen & 6 & 9 & 5（3） & \textbf{2}（2） & 0/0 & 2 & \textbf{2}（2） & \textbf{2}（2） \\
Sod & 0 & 1 & 5（0） & \textbf{0}（0） & 1/0 & 0 & \textbf{0}（0） & \textbf{0}（0） \\
爆炸波 & 0 & 0 & 10（0） & \textbf{0}（0） & 0/0 & 0 & \textbf{0}（0） & \textbf{0}（0） \\
激波–气泡 & 22 & 31 & 10（10） & \textbf{8}（8） & 1/0 & 2 & \textbf{3}（2） & \textbf{1}（1） \\
激波–涡 & 3 & 3 & 6（1） & \textbf{5}（1） & 0/0 & 3 & \textbf{3}（2） & \textbf{1}（1） \\
\bottomrule\end{tabular}}\end{table}
\begin{table}[H]\centering\small\caption{周长判据边界的噪声 IoU（与干净边界相比），三个版本。}\begin{tabular}{llccc}\toprule 算例 & 核 & 周长判据：第 1 版（2 次） & 第 2 版（$4(r/3h)^2$） & 第 3 版（$w\le r/6$）\\\midrule
双环流 & (1.51, 0.38) & 0.58 & 0.92 & 0.92\\
双环流 & (0.75, 0.73) & 0.87 & 0.93 & 0.93\\
Lamb–Oseen & (-3.04, -0.04) & 0.94 & 0.80 & 0.94\\
Lamb–Oseen & (3.03, 0.07) & 0.95 & 0.97 & 0.95\\
激波–涡 & (0.70, 0.50) & 0.58 & 0.78 & 0.78\\
激波–涡 & (0.84, 0.50) & — & 0.62 & —\\
激波–气泡 & (1.00, 0.04) & 0.58 & 0.44 & 0.58\\
\bottomrule\end{tabular}\end{table}
\begin{itemize}
\item 涡数与第 2 版相同：双环流 2/2、Lamb–Oseen 2/2、Sod 0/0、爆炸波 0/0、激波–气泡 3/1、激波–涡 3/1。
\item 第 3 版恢复了第 2 版中变差的两个边界：Lamb–Oseen 从 0.80 回到 0.94，激波–气泡小涡从 0.44 回到 0.58。
\item 激波–涡第二个涡在噪声下未找到边界。
\item 凸性判据在 Lamb–Oseen 上从 0.98 降到 0.86，因为小涡现在只平滑 2 次。
\end{itemize}

\section{结论与修订建议}
\begin{enumerate}
\item \keyp{分工。}
  \begin{itemize}[leftmargin=1.2em]
  \item 边界：由 \AFE 色散门导出的物质周长判据。
  \item 强度：相干环路上的环量，以外侧环带为参照。
  \item $\varphi$：只提供候选等值线与核的位置。
  \item 合并树 $P$：作为不依赖集合参照的核检测场的候选，前提是簇的环境在观测区域内。
  \end{itemize}
\item \keyp{参数。}$n=2+60\hat\sigma/h$，且 $w(n)\le r_{\rm eq}/6$；$J_{\max}=0.35$；报告半径–$\delta$ 曲线。
\item \keyp{理论稿要改的地方。}
  \begin{itemize}[leftmargin=1.2em]
  \item 写明逐点转角场（LAVD、$\varphi$）在涡片上没有良定义的极限；
  \item 强度改用环路积分；
  \item 写明同步语言与 LAVD 的对应：强度是频率失谐，层级给出局部参照。
  \end{itemize}
\item \keyp{尚未完成。}外侧环路的选取（激波–涡、激波–气泡大多找不到）；把环量强度接进流水线；合并树在噪声下的稳健性。
\end{enumerate}

\section*{参考文献}
\small
[1] A. P. Millán, J. J. Torres, G. Bianconi. Explosive higher-order Kuramoto dynamics on simplicial complexes. Phys. Rev. Lett. 124, 218301 (2020).\par\noindent [2] R. Ghorbanchian, J. G. Restrepo, J. J. Torres, G. Bianconi. Higher-order simplicial synchronization of coupled topological signals. Commun. Phys. 4, 120 (2021).\par\noindent [3] L. Calmon, J. G. Restrepo, J. J. Torres, G. Bianconi. Dirac synchronization is rhythmic and explosive. Commun. Phys. 5, 253 (2022).\par\noindent [4] L. M. Pecora, F. Sorrentino, A. M. Hagerstrom, T. E. Murphy, R. Roy. Cluster synchronization and isolated desynchronization in complex networks with symmetries. Nat. Commun. 5, 4079 (2014).\par\noindent [5] F. P. Kemeth, S. W. Haugland, L. Schmidt, I. G. Kevrekidis, K. Krischer. A classification scheme for chimera states. Chaos 26, 094815 (2016).\par\noindent [6] M. Budišić, I. Mezić. Geometry of the ergodic quotient reveals coherent structures in flows. Physica D 241, 1255–1269 (2012).\par\noindent [7] N. Tarshish, R. Abernathey, C. Zhang, C. O. Dufour, I. Frenger, S. M. Griffies. Identifying Lagrangian coherent vortices in a mesoscale ocean model. Ocean Modelling 130, 15–28 (2018).\par\noindent [8] A. Hadjighasem, D. Karrasch, H. Teramoto, G. Haller. Spectral-clustering approach to Lagrangian vortex detection. Phys. Rev. E 93, 063107 (2016).

\normalsize
\section*{复现}
\begin{itemize}
\item \texttt{tcx.py}：\texttt{[rev3]} 段改了 \texttt{n\_smooth} 的上限。
\item \texttt{tcx\_run.py}：$J_{\max}=0.35$。
\item 新增脚本：
  \begin{itemize}[leftmargin=1.2em]
  \item \texttt{sweep.py}、\texttt{sweep\_sum.py}：平滑与 $\delta$ 扫描；
  \item \texttt{neff.py}：单步占比与参与比；
  \item \texttt{sync\_lib.py}、\texttt{sync\_main.py}、\texttt{sync2.py}：合并树与集合检验；
  \item \texttt{sheet.py}：涡片检验；
  \item \texttt{circ\_strength.py}：环量强度；
  \item \texttt{afe\_tables3.py}、\texttt{afe\_tables4.py}：本报告的表。
  \end{itemize}
\item 六算例重跑：\texttt{run\_all\_rev3.sh}。
\end{itemize}
\end{document}
